Для апробации распознания текста проведён тест на десяти изображениях, содержащих только текст. Рассматривались шесть конфигураций обработки: базовый EasyOCR без дополнительных этапов (в таблице~\ref{tab:ocr_compare} обозначен за OCR), OCR с адаптивной предобработкой изображения (в таблице~\ref{tab:ocr_compare} обозначен за OCR+пре), OCR с нормализацией текста при помощи языковой модели (в таблице~\ref{tab:ocr_compare} обозначен за OCR+норм), OCR с эвристической нормализацией (в таблице~\ref{tab:ocr_compare} обозначен за OCR+нормЭ), а также две комбинированные схемы, включающие предобработку и последующую нормализацию текста (в таблице~\ref{tab:ocr_compare} обозначен за OCR+пред+нормЭ и OCR+пред+норм).

\begin{table*}[t]
	\centering
	\small
	\setlength{\tabcolsep}{4pt}
	\renewcommand{\arraystretch}{1.15}
	\resizebox{\textwidth}{!}{
		\begin{tabular}{c c c c c c c}
	\hline
	Тест & OCR & OCR+пред & OCR+норм & OCR+предЭ & OCR+пред+нормЭ & OCR+пред+норм \\
	\hline
	1  & 0,1603 & 0,4737 & 0,5465 & 0,1603 & 0,4503 & 0,9394 \\
	2  & 0,1742 & 0,5023 & 0,4838 & 0,1559 & 0,4909 & 0,8909 \\
	3  & 0,2151 & 0,0154 & 0,3588 & 0,2000 & 0,0310 & 0,4113 \\
	4  & 0,8734 & 0,8518 & 0,9240 & 0,8734 & 0,8500 & 0,9500 \\
	5  & 0,6488 & 0,7610 & 0,8750 & 0,5780 & 0,7121 & 0,9212 \\
	6  & 0,7091 & 0,6139 & 0,9643 & 0,7529 & 0,6450 & 0,8060 \\
	7  & 0,2100 & 0,3986 & 0,3720 & 0,1995 & 0,4203 & 0,8711 \\
	8  & 0,3656 & 0,0000 & 0,7923 & 0,3771 & 0,0000 & 0,0538 \\
	9  & 0,7210 & 0,7812 & 0,8785 & 0,7494 & 0,8031 & 0,8976 \\
	10 & 0,0636 & 0,3476 & 0,4295 & 0,0266 & 0,3523 & 0,9516 \\
	\hline
	Среднее & 0,4141 & 0,4746 & 0,6625 & 0,4073 & 0,4755 & 0,7693 \\
	Ст. откл. & 0,2936 & 0,2975 & 0,2460 & 0,3053 & 0,2935 & 0,2983 \\
	\hline
\end{tabular}
	}
	\caption{Результаты распознавания текста при различных конфигурациях OCR}
	\label{tab:ocr_compare}
\end{table*}

Анализ результатов показывает, что наилучшее качество распознавания обеспечивает комбинированная схема \textit{OCR+пре+норм}. Для нее получено максимальное среднее значение~--- 0,7693, что свидетельствует о высокой эффективности совместного применения адаптивной предобработки изображения и последующей нормализации текста языковой моделью.
Конфигурация \textit{OCR+норм} также демонстрирует высокий результат и наименьшее стандартное отклонение среди всех рассмотренных вариантов, что указывает на стабильность постобработки с помощью языковой модели на разных изображениях.
Использование только предобработки в схеме \textit{OCR+пре} дает лишь умеренное улучшение по сравнению с базовым OCR, а эвристическая нормализация уступает LLM-подходу как по среднему качеству, так и по общей эффективности. Комбинация предобработки с эвристиками также показывает ограниченный прирост и не достигает результатов схем с языковой моделью.

Таким образом, ключевой вклад в повышение качества распознавания вносит именно этап постобработки текста, особенно с использованием YandexGPT 5 Lite. Наиболее эффективной по совокупности результатов является конфигурация \textit{OCR+пре+норм}, что показывает, что предобработки изображения и постобработка текста сильно улучшают качества работы OCR в среднем на $35\%$.